% =============================================================================
% 1.3 Gestion documentaire et systèmes d'information dans le BTP
% =============================================================================
\section{Gestion documentaire et systèmes d'information dans le BTP}
\label{sec:ch1-ged}

\subsection{Évolution des systèmes de gestion documentaire}

Les systèmes de gestion électronique des documents ont connu une évolution
significative depuis les années 1990. Le tableau~\ref{tab:generations-ged}
en retrace les grandes générations.

\begin{btptable}{Générations de systèmes de gestion documentaire}
\label{tab:generations-ged}
\begin{tabularx}{\linewidth}{@{}L{2.5cm} L{1.9cm} Y Y@{}}
\toprule
\rowcolor{btppale}
\btpth{Génération} & \btpth{Période} & \btpth{Caractéristiques} & \btpth{Limites} \\
\midrule
Première & 1990--2000 & Archivage électronique, indexation manuelle. & Recherche limitée, aucune analyse sémantique. \\
Deuxième & 2000--2010 & Circuits de validation, métadonnées structurées. & Rigidité des schémas, coût d'implémentation. \\
Troisième & 2010--2020 & Hébergement infonuagique, collaboration, recherche plein texte. & Exploitation superficielle du contenu. \\
Quatrième & Depuis 2020 & Analyse sémantique, extraction et génération assistées. & Maturité en cours de consolidation. \\
\bottomrule
\end{tabularx}
\end{btptable}

La quatrième génération, qui intègre des capacités d'intelligence artificielle,
est particulièrement prometteuse pour le domaine des appels d'offres, car elle
permet non seulement de rechercher des documents, mais aussi d'en extraire des
connaissances et de générer de nouveaux contenus. C'est dans cette génération
que se situe notre contribution.

\subsection{Standards et normes applicables}

Plusieurs référentiels encadrent la gestion documentaire dans le secteur~:

\begin{itemize}
  \item \btpterm{ISO 15489} — gestion des documents d'activité~: définit les
        principes de création, de capture et de gestion des documents engageants~;
  \item \btpterm{ISO 30301} — systèmes de management des documents d'activité~:
        cadre organisationnel et exigences d'audit~;
  \item \btpterm{Modélisation des données du bâtiment} (\anglais{Building
        Information Modeling})~: structuration des données techniques d'un
        ouvrage sur son cycle de vie.
\end{itemize}

Ces référentiels partagent une exigence directement pertinente pour notre
travail~: la \btpterm{traçabilité}. Un document engageant doit pouvoir être
rattaché à son auteur, à sa version et à ses sources. Cette exigence normative
converge avec la contrainte technique que nous imposerons au système de
génération~: toute affirmation produite doit être citable.

\subsection{Les limites des solutions actuelles}

Les solutions de gestion documentaire et de gestion des appels d'offres
disponibles sur le marché présentent quatre limitations structurelles.

\begin{enumerate}
  \item \btpterm{Recherche sémantique limitée.} La plupart des systèmes
        fonctionnent sur la base de mots-clés, sans compréhension du sens ni du
        contexte. Une recherche de « pénalités de retard » ne retrouve pas un
        passage évoquant les « sanctions pour dépassement de délai ».

  \item \btpterm{Absence d'exploitation des relations.} Les liens entre
        documents, projets, clients, certifications et compétences ne sont pas
        explicitement modélisés. L'information existe, mais les chemins qui la
        relient restent implicites.

  \item \btpterm{Capitalisation passive.} Les connaissances restent enfermées
        dans les documents, sans extraction ni réutilisation active. Le système
        restitue des fichiers, non des réponses.

  \item \btpterm{Absence de génération.} Les systèmes ne produisent pas de
        nouveaux contenus à partir du patrimoine existant. La rédaction demeure
        intégralement humaine, y compris pour des sections largement
        réutilisables d'une offre à l'autre.
\end{enumerate}

\begin{btpwarn}[Ce que l'automatisation ne doit pas faire disparaître]
Ces limites pourraient laisser croire qu'une automatisation intégrale est
souhaitable. Ce serait une erreur d'analyse. Le caractère engageant de l'offre
impose que la responsabilité reste humaine. L'objectif n'est donc pas de
supprimer l'intervention de l'ingénieur, mais de déplacer son effort~: de la
recherche et de la transcription vers la vérification et l'arbitrage.
\end{btpwarn}
